\section{Related Work}\label{sec:related}
\subsection{Few-Shot Industrial Anomaly Detection}
Classical industrial anomaly detection (AD) typically assumes access to normal-only training data and identifies deviations at test time. Patch distribution, self-supervised, and reconstruction-discriminative approaches such as PaDiM, CutPaste, and DRAEM provide important context for this literature~\cite{padim2021,cutpaste2021,draem2021}. PatchCore~\cite{roth2022towards} established memory-bank patch retrieval as a robust baseline by storing representative normal features and scoring test patches using nearest-neighbor distance. AnomalyDINO~\cite{damm2025anomalydino} extends this paradigm to few-shot inspection with frozen DINOv2 representations. The few-shot AD literature has further introduced registration-based adaptation (RegAD~\cite{regad2022}), graph-structured memory (GraphCore~\cite{graphcore2023}), fast feature reconstruction (FastRecon~\cite{fastrecon2023}), vision-language zero- and few-shot scoring (WinCLIP~\cite{winclip2023}; CLIP-SAM collaboration~\cite{clipsam2025} and cross-modal prompt regularization~\cite{uniadprompt2026} continue this direction), and in-context residual learning (InCTRL~\cite{inctrl2024}). While these methods are effective rankers, the reference state can remain nontrivial at inference time, and the raw score scale does not represent a calibrated probability. Calibration under distribution shift is also known to degrade for deep models generally~\cite{ovadia2019trust}. Our work is complementary: it does not compete on ranking but instead examines how few-shot detectors should expose decision reliability.

\subsection{Frozen-Feature Subspace Detectors}
Subspace approaches replace explicit patch retrieval with a compact normal subspace. SubspaceAD~\cite{subspacead2026} demonstrates that frozen DINOv2 patch features and PCA or subspace residuals are already effective for few-shot anomaly ranking, establishing the ranking mechanism as settled prior art. Our contribution begins after the residual is computed, focusing on how a low-storage subspace detector should expose reliability under low-shot and distribution shift conditions, as well as the guarantees its alarms can provide.

\subsection{Calibration and Robustness in Few-Shot AD}
Calibration evaluates whether predicted confidence aligns with empirical correctness~\cite{platt1999,guo2017calibration}. Standard post-hoc calibrators include Platt scaling~\cite{platt1999}, temperature scaling~\cite{guo2017calibration}, histogram binning~\cite{zadrozny2001obtaining}, and isotonic regression~\cite{zadrozny2002transforming}; all are compared under a shared label-free protocol. In anomaly detection, calibration is challenging due to sparse labels, open-ended anomaly types, and varying score distributions across classes. Reliability analyses of DINOv2-based few-shot AD report poor expected calibration error (ECE) and significant FGSM fragility~\cite{khan2025calibration}. These findings motivated us to evaluate ECE, Brier score, and negative log-likelihood (NLL) separately from AUROC and average precision (AP), and to report adversarial results, including FGSM-style perturbations~\cite{goodfellow2015}, only as fragility diagnostics rather than as evidence of improved robustness.

\subsection{Conformal Prediction and Risk-Controlled Calibration}
Conformal prediction transforms nonconformity scores into p-values under exchangeability assumptions~\cite{vovk2005,shafer2008tutorial,angelopoulos2021gentle}. This approach is particularly suitable for normal-only few-shot AD, as support normal images can define a reference distribution without requiring real anomalies. Conformal anomaly detection was already well established: leave-one-out, bootstrap, and cross-conformal anomaly detectors had been directly analyzed~\cite{hennhofer2024loio}; weighted conformal detection in low-data regimes had examined resolution and effective-sample-size effects~\cite{hennhofer2026weighted}; general nonconformity toolkits had been developed~\cite{hennhofer2026nonconform}; and few-shot conformal prediction with auxiliary tasks had addressed small target calibration sets~\cite{fisch2021fewshot}. Conformal AD, LOIO calibration, auxiliary-task transfer, and weighted conformal p-values are therefore established mechanisms.

Risk-controlling calibration also preceded our work. Learn then Test cast finite-sample calibration as a multiple-testing problem over candidate configurations~\cite{angelopoulos2021learntest}, while conformal risk control selected a parameter to control the expectation of a monotone loss~\cite{angelopoulos2024crc}. In anomaly detection, PAC-Wrap provided semi-supervised PAC bounds for false-positive and false-negative rates around a base detector~\cite{li2022pacwrap}. These works established the general principles of candidate testing, held-out risk assessment, and one-sided error bounds; we do not claim those statistical primitives as new. CRESS instead instantiates them for normal-only cross-category thresholding by assigning disjoint source categories to reference, proposal, and certification roles, distinguishing fixed-archive image risk from new-category risk, and auditing whether the available number of independent categories can support a nonzero threshold.

Structure-aware conformal adaptation had also been studied outside anomaly detection. Lam and Anh proposed HeAD-CP, which adapted graph-score diffusion using label-free local-homophily estimates and showed that uniform diffusion could impair prediction-set efficiency on heterophilic graphs while preserving marginal coverage~\cite{lam2026headcp}. This result supports the broader need to audit reliability transformations under structural shift, but its setting concerns GNN node-classification sets rather than few-shot industrial anomaly alarms or cross-category source certification.

The present contribution begins at the few-shot resolution floor $1/(k{+}1)$ below which no deterministic target-only $k$-score conformal alarm can operate, and at the independent-category requirement of a simultaneous source certificate. This work establishes the algebraic floor and the Hoeffding-specific and distribution-free all-zero category-feasibility bounds, empirically audits corruption-shift deviations, and evaluates source-assisted candidate thresholds without target anomaly labels. CRESS does not introduce a new nonconformity score or generic confidence bound; it specifies how source categories are separated into reference, proposal, and certification roles and audits whether the resulting source certificate can be nonvacuous at the available category count.
